\section{Separate the tap walker from the output window}
The scalar path uses \texttt{t0} for its exact accumulator, \texttt{t1} and
\texttt{t2} for private input/coefficient walkers, and \texttt{t3} for remaining taps.
\texttt{a0} identifies the outer output's input-window base; \texttt{a2} identifies
its destination sample; \texttt{a3} counts remaining outputs. Both private walkers
restart each output. \texttt{v0} retains the original count and \texttt{v1} accumulates
clipped outputs. A word-sized sample does not justify a word-sized address operation.

The vector path retains widened coefficient words in \texttt{w3} and \texttt{w4}.
It reconstructs input words in \texttt{w5} and \texttt{w6} for every window. Dot
products and their reduction use \texttt{w7}, \texttt{w8}, \texttt{t0} and \texttt{t1}.
These roles explain this sequence; final free-writing grading does not require them.
The separately retained \texttt{fir\_scalar} is a guided comparison routine, not an
additional public entry required of an equivalent submission.

\section{Prepare a fresh construction observation}
The companion \texttt{packages/021-fixed-point-fir/starter/main.asm} supplies both function bodies
with marked regions. The \texttt{input} array in each recipe includes all seven
history samples and the requested new samples. Thus output count is its length minus
seven. The \texttt{taps} array always has eight signed coefficient codes in oldest-first
order. Preserve the 192-byte input and output arenas and 64-byte coefficient arena.
The preparation helper fills the output arena with hexadecimal a5, uses the listed
arrays and zero padding, and preserves guards and caller metadata.

Start a fresh Debug execution and stop before \texttt{filter\_ready}. The caller has
not yet stored the result words. The expected output prefix is followed by a5 bytes
where a partial program has not stored output. Vector lanes in the recipes are raw
unsigned bit patterns at the stated width; signed scalar expectations use full signed
register values. The recipe guide supplies the exact source-preparation command.

Private construction checks repeat each recipe through this caller and through an
isolated call with input, coefficient and destination byte offsets three, five and
seven. These are separate observation contexts. A provisional return from unfinished
code is construction evidence, not a completed filtering result. Numerical models and automated light-theme captures verify the distinct boundaries
shown in the figures.

The student package includes public construction scripts in
\texttt{checks/construction}. Run the named script against the cumulative program
before proceeding. Step forty-two also has a completed-program check for its two
rounded-clipping recipes. These guided scripts may name the illustrated temporaries;
final free-writing checks enforce the public interface and ownership instead.
